\section{Grok Bot、Enterpriseとテンプレートを訴求}

Elon MuskはGrok Botのテンプレート機能とEnterprise向け提供を取り上げ、調達交渉を想定したHaggle Botなどの例を紹介した。最近発表されたGrok Bot for Enterpriseの利用イメージを具体化する話題として、汎用チャットではなく特定業務の交渉・処理手順をテンプレート化したエージェントへ落とし込む方向性が注目された。

Grok収集では、企業利用とテンプレートの双方についてMusk自身の窓内投稿が確認されている。今回の観測では機能の網羅的な仕様や価格よりも、調達のような定型業務を「Botの型」として再利用し、削減額などの成果を訴求する製品ポジショニングが前面に出た。

\dailyxsource{Elon Musk (@elonmusk)}{https://x.com/elonmusk/status/2095988111222227130}
